% GENERATED FILE. Edit canonical lessons, then run npm run generate:books.
% canonical-source-hash: fnv1a64:5380fad648ec89a7
% canonical-lessons: MR-C452-dhakka, MR-C452-sahara, MR-C452-dhasakne, MR-C452-bitharne, MR-C452-habakne

\chapter{A Sudden Fright}
\label{ch:mr-a-sudden-fright}
\hlchaptermodality{452}

\begin{chapteropening}
\textbf{By the end of this chapter:} \emph{I can describe a shock and a sudden fright.}

Complete the last lesson of chapter 452: I can describe a shock and a sudden fright.
\end{chapteropening}

\section[\texorpdfstring{\mr{धक्का} (dhakkā)}{धक्का (dhakkā)}]{\mr{धक्का} (dhakkā) — a shock; a push}

\label{lesson:MR-C452-dhakka}



\begin{quote}
Before the new one: say the Marathi for displeased, then the Marathi for to get fed up.
\end{quote}

\subsection*{You'll want to know: \mr{धक्का}}
\textbf{\mr{धक्का}} — \emph{dhakkā} — \textquotedblleft{}a shock; a push\textquotedblright{}.

\begin{quote}
\mr{ती} \mr{बातमी} \mr{ऐकून} \mr{आम्हाला} \mr{मोठा} \mr{धक्का} \mr{बसला}.
\end{quote}

\begin{quote}
\emph{tī bātmī aikūn āmhālā moṭhā dhakkā baslā.}
\end{quote}

\begin{quote}
Hearing that news gave us a big shock.
\end{quote}

Masculine.

\mr{धक्का} \mr{बसणे}, to get a shock, with the dative; \mr{धक्का} \mr{देणे}, to push.

\begin{cousinweb}
\textbf{\mr{धक्का}} (\emph{dhakkā}) is inherited from the expressive root \emph{*dhakk}, \textquotedblleft{}push, strike\textquotedblright{} (CDIAL 6701), under which CDIAL lists Marathi \emph{dhakā}, \textquotedblleft{}push\textquotedblright{}.
\end{cousinweb}

\begin{grammarlens}[title={What you've built}]
One more word for describing what you lived and felt.
\end{grammarlens}

\subsection*{Your turn}
\begin{itemize}
\raggedright
  \item \emph{Say it:} \emph{dhakkā}
  \item \emph{Say it:} \emph{dhakkā}, once more
  \item \emph{Say it:} \emph{tī bātmī aikūn āmhālā moṭhā dhakkā baslā.} — Hearing that news gave us a big shock.
  \item \emph{Recall:} say the Marathi for displeased, then the Marathi for to get fed up, then say \emph{dhakkā} again
\end{itemize}

\begin{tcolorbox}[breakable,colback=teal!4,colframe=teal!35!black,title={Before you move on}]
What does \mr{धक्का} mean? (\textquotedblleft{}A shock; a push\textquotedblright{}.) Say it once more.
\end{tcolorbox}
\section[\texorpdfstring{\mr{शहारा} (śahārā)}{शहारा (śahārā)}]{\mr{शहारा} (śahārā) — a shiver; goosebumps}

\label{lesson:MR-C452-sahara}



\begin{quote}
Before the new one: say the Marathi for to get fed up, then the Marathi for a shock.
\end{quote}

\subsection*{You'll want to know: \mr{शहारा}}
\textbf{\mr{शहारा}} — \emph{śahārā} — \textquotedblleft{}a shiver; goosebumps\textquotedblright{}.

\begin{quote}
\mr{ती} \mr{गोष्ट} \mr{ऐकून} \mr{मला} \mr{शहारा} \mr{आला}.
\end{quote}

\begin{quote}
\emph{tī goṣṭa aikūn malā śahārā ālā.}
\end{quote}

\begin{quote}
Hearing that story gave me a shiver.
\end{quote}

Masculine.

\mr{शहारा} \mr{येणे}, to shiver, with the dative. The plural \mr{शहारे}, goosebumps all over.

\begin{cousinweb}
\textbf{\mr{शहारा}} (\emph{śahārā}) is a Marathi noun whose origin is not settled.
\end{cousinweb}

\begin{grammarlens}[title={What you've built}]
One more word for describing what you lived and felt.
\end{grammarlens}

\subsection*{Your turn}
\begin{itemize}
\raggedright
  \item \emph{Say it:} \emph{śahārā}
  \item \emph{Say it:} \emph{śahārā}, once more
  \item \emph{Say it:} \emph{tī goṣṭa aikūn malā śahārā ālā.} — Hearing that story gave me a shiver.
  \item \emph{Recall:} say the Marathi for to get fed up, then the Marathi for a shock, then say \emph{śahārā} again
\end{itemize}

\begin{tcolorbox}[breakable,colback=teal!4,colframe=teal!35!black,title={Before you move on}]
What does \mr{शहारा} mean? (\textquotedblleft{}A shiver; goosebumps\textquotedblright{}.) Say it once more.
\end{tcolorbox}
\section[\texorpdfstring{\mr{धसकणे} (dhasakṇe)}{धसकणे (dhasakṇe)}]{\mr{धसकणे} (dhasakṇe) — to take fright}

\label{lesson:MR-C452-dhasakne}



\begin{quote}
Before the new one: say the Marathi for a shock, then the Marathi for a shiver.
\end{quote}

\subsection*{You'll want to know: \mr{धसकणे}}
\textbf{\mr{धसकणे}} — \emph{dhasakṇe} — \textquotedblleft{}to take fright\textquotedblright{}.

\begin{quote}
\mr{दार} \mr{अचानक} \mr{वाजलं} \mr{आणि} \mr{आजी} \mr{धसकली}.
\end{quote}

\begin{quote}
\emph{dār acānak vājlaṁ āṇi ājī dhasaklī.}
\end{quote}

\begin{quote}
The door banged suddenly and Grandmother took fright.
\end{quote}

Past (he …): \mr{तो} \mr{धसकला}, \emph{to dhasaklā}.

\mr{धसका}, a sudden fright; \mr{धसका} \mr{घेणे}, to take something badly to heart.

\begin{cousinweb}
\textbf{\mr{धसकणे}} (\emph{dhasakṇe}) is a Marathi verb whose origin is not settled.
\end{cousinweb}

\begin{grammarlens}[title={What you've built}]
One more word for describing what you lived and felt.
\end{grammarlens}

\subsection*{Your turn}
\begin{itemize}
\raggedright
  \item \emph{Say it:} \emph{dhasakṇe}
  \item \emph{Say it:} \emph{dhasakṇe}, once more
  \item \emph{Say it:} \emph{dār acānak vājlaṁ āṇi ājī dhasaklī.} — The door banged suddenly and Grandmother took fright.
  \item \emph{Recall:} say the Marathi for a shock, then the Marathi for a shiver, then say \emph{dhasakṇe} again
\end{itemize}

\begin{tcolorbox}[breakable,colback=teal!4,colframe=teal!35!black,title={Before you move on}]
What does \mr{धसकणे} mean? (\textquotedblleft{}To take fright\textquotedblright{}.) Say it once more.
\end{tcolorbox}
\section[\texorpdfstring{\mr{बिथरणे} (bitharṇe)}{बिथरणे (bitharṇe)}]{\mr{बिथरणे} (bitharṇe) — to panic; to run wild}

\label{lesson:MR-C452-bitharne}



\begin{quote}
Before the new one: say the Marathi for a shiver, then the Marathi for to take fright.
\end{quote}

\subsection*{You'll want to know: \mr{बिथरणे}}
\textbf{\mr{बिथरणे}} — \emph{bitharṇe} — \textquotedblleft{}to panic; to run wild\textquotedblright{}.

\begin{quote}
\mr{फटाक्याचा} \mr{आवाज} \mr{ऐकून} \mr{घोडा} \mr{बिथरला}.
\end{quote}

\begin{quote}
\emph{phaṭākyācā āvāj aikūn ghoḍā bitharlā.}
\end{quote}

\begin{quote}
The horse panicked at the sound of the firecracker.
\end{quote}

Past (he …): \mr{तो} \mr{बिथरला}, \emph{to bitharlā}.

Of animals and people who lose their heads in fright: \mr{घोडा} \mr{बिथरला}, the horse bolted.

\begin{cousinweb}
\textbf{\mr{बिथरणे}} (\emph{bitharṇe}) is inherited by CDIAL's account: it lists Marathi \emph{bitharṇẽ}, \textquotedblleft{}to misconduct oneself\textquotedblright{}, beside \emph{itharṇẽ}, \textquotedblleft{}to run wild\textquotedblright{}, under Sanskrit \emph{itvara}, \textquotedblleft{}walking; mean, cruel\textquotedblright{} (CDIAL 1566). How the \emph{b-} arose is not explained there, so the link is not certain.
\end{cousinweb}

\begin{grammarlens}[title={What you've built}]
One more word for describing what you lived and felt.
\end{grammarlens}

\subsection*{Your turn}
\begin{itemize}
\raggedright
  \item \emph{Say it:} \emph{bitharṇe}
  \item \emph{Say it:} \emph{bitharṇe}, once more
  \item \emph{Say it:} \emph{phaṭākyācā āvāj aikūn ghoḍā bitharlā.} — The horse panicked at the sound of the firecracker.
  \item \emph{Recall:} say the Marathi for a shiver, then the Marathi for to take fright, then say \emph{bitharṇe} again
\end{itemize}

\begin{tcolorbox}[breakable,colback=teal!4,colframe=teal!35!black,title={Before you move on}]
What does \mr{बिथरणे} mean? (\textquotedblleft{}To panic; to run wild\textquotedblright{}.) Say it once more.
\end{tcolorbox}
\section[\texorpdfstring{\mr{हबकणे} (habakṇe)}{हबकणे (habakṇe)}]{\mr{हबकणे} (habakṇe) — to be stunned}

\label{lesson:MR-C452-habakne}



\begin{quote}
Before the new one: say the Marathi for to take fright, then the Marathi for to panic.
\end{quote}

\subsection*{You'll want to know: \mr{हबकणे}}
\textbf{\mr{हबकणे}} — \emph{habakṇe} — \textquotedblleft{}to be stunned\textquotedblright{}.

\begin{quote}
\mr{तो} \mr{अपघात} \mr{पाहून} \mr{सगळे} \mr{हबकले}.
\end{quote}

\begin{quote}
\emph{to apghāt pāhūn sagḷe habakle.}
\end{quote}

\begin{quote}
Everyone was stunned to see that accident.
\end{quote}

Past (he …): \mr{तो} \mr{हबकला}, \emph{to habaklā}.

\mr{हबकून} \mr{जाणे}, to be struck dumb with shock.

\begin{cousinweb}
\textbf{\mr{हबकणे}} (\emph{habakṇe}) is a Marathi verb whose origin is not settled.
\end{cousinweb}

\begin{grammarlens}[title={What you've built}]
One more word for describing what you lived and felt.
\end{grammarlens}

\subsection*{Your turn}
\begin{itemize}
\raggedright
  \item \emph{Say it:} \emph{habakṇe}
  \item \emph{Say it:} \emph{habakṇe}, once more
  \item \emph{Say it:} \emph{to apghāt pāhūn sagḷe habakle.} — Everyone was stunned to see that accident.
  \item \emph{Recall:} say the Marathi for to take fright, then the Marathi for to panic, then say \emph{habakṇe} again
\end{itemize}

\begin{tcolorbox}[breakable,colback=teal!4,colframe=teal!35!black,title={Before you move on}]
What does \mr{हबकणे} mean? (\textquotedblleft{}To be stunned\textquotedblright{}.) Say it once more.
\end{tcolorbox}
